% Introduction section. Paste into your own thesis, or \input it.
%
% TODO (CGPS): this section deliberately carries NO exact figures from the Copenhagen General
% Population Study -- cohort size, prevalence split, follow-up, event counts. Data access is
% pending; the numbers go in once the cohort actually available to this project is known.
% The qualitative claims below are sourced to fuchs2023cgps and hold without them.
%
% This is the thesis-level problem statement: silent disease -> why prevention needs a measurable
% signal in the well -> plaque is focal, not uniform, and geometry explains why -> that link is
% documented only in trees disease has already reshaped -> the untouched question is the topology
% of trees before that reshaping -> what this thesis does about it -> what bounds the answer.
%
% Rewritten 2026-09-07, replacing an earlier draft of this section in full (not a polish pass).
% STRUCTURE (after Aida Jimenez's MSc thesis ch. 1, Former_students_work/): this section is the
% funnel only -- stake, mechanism, evidence, flaw, gap -- and stops at the gap. The aim, the
% research questions and the limits that bound them live in Section~\ref{sec:aims}.
%
% RELATION TO THE STATE OF THE ART. Per docs_thesis/state_of_the_art_plan.md, the state of the art
% opens on the same scale argument but carries the evidence matrix and surveys the field study by
% study. This section states the argument with the few numbers that carry it; that chapter proves
% it in full. Do not expand this into a survey.
%
% Verification status of the sources used here: griffo2026wss and shen2026geometry are full-text
% read (docs_thesis/papers/). The rest are abstract-level, verified against Crossref/Europe PMC
% metadata in docs_thesis/find_research.md.

\section{Framing}
\label{sec:framing}

Coronary artery disease is characteristically silent: it advances for years before it produces a
symptom, and for many patients the first sign of it is the heart attack itself
\cite{knuuti2019esc}.
Screening adults who have no known heart disease with \textbf{coronary computed tomography
angiography (CCTA)} shows how far that silence extends.
Among people with no reason before the scan to suspect anything, a large share are found to carry
coronary atherosclerosis already, most of it not yet obstructing blood flow and some of it severe
enough to do so \cite{fuchs2023cgps}.

Treating the disease before that point means finding people while they still feel well, which means
finding something measurable in a population with no complaint to prompt the scan.
Cholesterol level and blood pressure are already measured this way, but neither predicts where in a
given patient's arteries the disease will actually appear.
Both bathe the whole coronary tree equally, while plaque does not \cite{chatzizisis2007ess}.
What is uniform is the exposure; what is not is the location, and the location is set by shape.

The mechanism runs through the vessel wall.
The endothelium, the single layer of cells lining the lumen, senses the drag the passing blood
exerts on it, the wall shear stress, and it responds differently depending on whether that drag is
steady or reverses through the cardiac cycle.
Low and reversing shear switches on the gene program that starts atherosclerosis: proteins that
recruit inflammatory cells to the wall, and a lining that lets lipid through it \cite{chatzizisis2007ess}.
Bends and branch points are where that pattern occurs, because they are where flow slows or briefly
reverses, and tracing the flow through five human coronary trees after death located plaque almost
exclusively there: on the outer wall of a vessel leaving a bifurcation, and along the inner wall of a
curved segment \cite{asakura1990flow}.
Shape does not merely correlate with where plaque forms; shape is the reason the same risk factors
produce disease in one place and not another.

That link is well documented, but only ever in trees a clinician already had reason to image.
In 313 patients split by angiographic stenosis, the left main bifurcation measured
$75.5^\circ \pm 14.7^\circ$ in 40 normal arteries, $81.2^\circ \pm 20.1^\circ$ in 62 with
non-significant disease, and $87.3^\circ \pm 18.8^\circ$ in 211 with a significant stenosis
\cite{juan2017bifurcation}.
The dependence goes beyond association.
Griffo et al.\ reconstructed 1078 coronary arteries from 748 patients pooled across three trials of
already-stenotic disease and trained a network to recover wall shear stress from vessel geometry
alone, with no flow simulation \cite{griffo2026wss}.
The network matched a full flow simulation at identifying which lesion in a patient's tree went on
to cause a heart attack: an area under the curve of 0.63 for both, statistically indistinguishable
($p=0.84$).
The figure itself is modest, but the equivalence is the finding.
Geometry alone carries essentially everything the simulation extracts, in vessels that were already
narrowed when they were measured.

That last detail is the problem.
An artery does not hold its shape while a lesion grows inside it.
Examining the left main artery in 136 hearts at autopsy, Glagov et al.\ found the vessel enlarges
outward to compensate, so the lumen does not measurably narrow until a lesion already occupies about
40\% of the vessel's cross-section \cite{glagov1987remodeling}.
Below that point the artery's outer wall is already growing around the plaque; above it, the lumen
falls away sharply.
Every geometry measured in a diseased artery is therefore measured after this compensation has
already begun, which means a cross-sectional study of diseased trees cannot cleanly tell a shape that
predisposed to the lesion apart from a shape the lesion has already produced.
The two cohorts above describe geometry that disease has already touched.

What has not been examined is the geometry of trees before that touch: the topology of a population
of coronary trees not selected for disease, described at a scale large enough to say what is
ordinary and what is not.
Coronary anatomy has been characterized as a population before, over more than 300 CCTA scans
chosen specifically for the absence of calcium and stenosis \cite{medranogracia2016atlas}, and
automatically in 281 patients stratified by cardiovascular risk factor
\cite{nannini2024characterization}.
Both stop at description.
Neither asks whether some of the trees in that normal range are more unusual than others, or whether
an unusual tree means anything for the person it belongs to.
A 2026 review of the field, written by several of the groups behind the mechanism above, reaches the
same limit from the literature side: work relating coronary anatomy to blood flow at population
scale remains rare, and almost all of it examines one bifurcation rather than the whole tree
\cite{shen2026geometry}.
Whether the shape of a disease-free coronary tree carries information about the person it belongs to,
and whether some shapes carry more of it than others, is open.
